\section{“到土里去”：理想主义的现实形态}

访谈后半段，主持人追问技术 vision。高继扬给出的愿景是：像培训员工一样培训机器人，通过几次示范和几次自我演练，让机器人在场景里稳定自主完成任务。为了实现这个体验，星海图需要基础模型、后训练工具和整机三件套。

\begin{figure}[H]
\centering
\includegraphics[width=0.94\textwidth]{go-to-soil-robotics.png}
\caption{“到土里去”：机器人创业为什么不浪漫。}
\end{figure}

\begin{knowledgebox}{读图：不浪漫不等于没有理想主义}
左侧软件/大模型应用更容易集中在模型、增长和传播；右侧机器人公司必须面对整机、供应链、客户现场、维护、安全和数据回流。中间的“土里”代表真实世界摩擦。机器人创业的理想主义不是口号，而是每天处理这些摩擦后仍然坚持目标。
\end{knowledgebox}

\subsection{如何识别画饼}

高继扬认为，无法在一个瞬间识别机器人公司是否画饼，需要看周期：团队一两年前说了什么，过去一年做了什么，一年后是否兑现。机器人行业里，画未来是必要的，因为员工、投资人、供应商和客户都需要相信未来；关键在于团队是否持续把未来描述变成现实。

\begin{importantbox}{画饼与战略叙事的区别}
描述未来不是问题；问题是未来是否被拆成阶段性结果、资源配置、客户价值和可验证交付。能兑现的叙事是战略，不能兑现的叙事才是空饼。
\end{importantbox}

\subsection{狼：进取心与韧性}

当被问到星海图像什么动物时，高继扬说最接近狼，但又强调每家公司都很狼性。这里的“狼”不是粗暴文化，而是面对客户交付和发布节点时的进取心、韧性和达成目标的劲头。具身智能行业的竞争强度，要求团队既有理想主义，又每天务实计算投入产出和长期价值。

\subsection{本章小结}

“到土里去”是本期最重要的产业隐喻。机器人公司必须在理想主义与务实之间找平衡：愿景要足够远，但每天都要解决硬件、供应链、客户、数据和组织问题。

